\chapter{Glossario}
\subsection*{A}
\begin{description}[leftmargin=2cm, labelindent=1cm]
    \item [Admin] utente con privilegi di gestione, in grado di creare e assegnare utenti e issue.
  \item[Affidabilità] misura della capacità del sistema di funzionare in modo corretto e continuativo nel tempo.
  \item[Angular] framework JavaScript per lo sviluppo di applicazioni web single-page.
  \item[Architettura] struttura logica e fisica che definisce l'organizzazione dei componenti software.
\end{description}

% B
\subsection*{B}
\begin{description}[leftmargin=2cm, labelindent=1cm]
  \item[Backend] parte dell'applicazione che gestisce la logica, il database e la comunicazione tra client e server.
  \item[Background tecnico] insieme delle conoscenze, competenze e tecnologie pregresse che costituiscono la base teorica e pratica di un team di sviluppo.
  \item[Boolean] tipo di dato che può assumere solo due valori: vero o falso.
  \item[Bug] errore o comportamento anomalo nel software che causa un malfunzionamento.
\end{description}

% C
\subsection*{C}
\begin{description}[leftmargin=2cm, labelindent=1cm]
  \item[Casi d'uso] rappresentazione delle interazioni tra un attore e il sistema per raggiungere un obiettivo specifico.
  \item[Chrome] browser web sviluppato da Google.
  \item[Click] azione dell'utente che interagisce con un elemento grafico.
  \item[Cloud computing] modello di erogazione di servizi informatici tramite Internet.
  \item[Code quality (Qualità del codice)] grado di leggibilità, modularità e manutenibilità del codice sorgente.
  \item[Commenti] annotazioni nel codice che descrivono funzionalità o scelte implementative.
  \item[Concorrenti] utenti o processi che operano simultaneamente sul sistema.
  \item[Containerizzato] modalità di distribuzione del software all'interno di contenitori isolati, ad esempio tramite Docker.
  \item[CSRF] vulnerabilità che consente ad un utente malintenzionato di eseguire azioni non autorizzate su un sito autenticato.
\end{description}

% D
\subsection*{D}
\begin{description}[leftmargin=2cm, labelindent=1cm]
  \item[Deadline] scadenza entro la quale una determinata attività o task deve essere completata.
  \item[Degrado] peggioramento delle prestazioni del sistema dovuto a carichi elevati.
  \item[Dipendenze] librerie o componenti esterni necessari al corretto funzionamento del software.
  \item[Docker] piattaforma per creare, distribuire ed eseguire container software.
\end{description}

% E
\subsection*{E}
\begin{description}[leftmargin=2cm, labelindent=1cm]
  \item[Edge] browser web sviluppato da Microsoft.
  \item[Efficienza] misura delle risorse necessarie per eseguire operazioni (tempo, memoria, CPU).
  \item[Esportare] operazione per convertire o salvare dati in formati condivisibili (CSV, PDF, ecc.).
  \item[Etichette] parole chiave che classificano o raggruppano issue e bug.
\end{description}

% F
\subsection*{F}
\begin{description}[leftmargin=2cm, labelindent=1cm]
  \item[File] entità che rappresenta un'informazione salvata in memoria.
  \item[File explorer] finestra di sistema usata per cercare e selezionare file locali.
  \item[Firefox] browser open source sviluppato da Mozilla.
  \item[Flusso] sequenza ordinata di azioni o eventi.
  \item[Formalismo] insieme delle regole che definiscono il modo corretto di rappresentare i casi d'uso.
  \item[Framework] struttura software riutilizzabile che facilita lo sviluppo applicativo.
  \item[Frontend] parte dell'applicazione visibile e interattiva con l'utente.
  \item[Funzionalità core] insieme delle funzioni principali del sistema.
\end{description}

% G
\subsection*{G}
\begin{description}[leftmargin=2cm, labelindent=1cm]
  \item[Gerarchia] relazione di subordinazione o dipendenza tra elementi o ruoli.
  \item[Git] sistema di controllo versione distribuito.
\end{description}

% H
\subsection*{H}
\begin{description}[leftmargin=2cm, labelindent=1cm]
  \item[Hashing] tecnica per trasformare dati (come password) in una stringa di lunghezza fissa non invertibile.
  \item[HEX / RRGGBB] codifica esadecimale utilizzata per rappresentare i colori.
  \item[Hibernate] framework Java per la gestione della persistenza dei dati.
\end{description}

% I
\subsection*{I}
\begin{description}[leftmargin=2cm, labelindent=1cm]
  \item[Icona] elemento grafico che rappresenta un'azione o funzione.
  \item[IDE] ambiente di sviluppo integrato che facilita la scrittura e il debug del codice.
  \item[Identificativo] codice univoco che distingue un elemento nel sistema.
  \item[Image] file grafico o rappresentazione visiva associata a un'issue.
  \item[Interfaccia] insieme degli elementi grafici e funzionali con cui l'utente interagisce.
  \item[Interazioni] comunicazioni o scambi di informazioni tra attori e sistema.
  \item[Issue] registrazione di un problema, richiesta o attività da gestire nel sistema.
\end{description}

% J
\subsection*{J}
\begin{description}[leftmargin=2cm, labelindent=1cm]
  \item[Jakarta EE] piattaforma Java per lo sviluppo di applicazioni enterprise.
  \item[Java] linguaggio di programmazione orientato agli oggetti.
  \item[JavaDoc] sistema di documentazione automatica per codice Java.
  \item[JSON Web Token (JWT)] standard aperto per la trasmissione sicura di informazioni tra client e server.
  \item[JS Doc] strumento per generare documentazione dal codice JavaScript.
\end{description}

% K (nessuna voce)

% L
\subsection*{L}
\begin{description}[leftmargin=2cm, labelindent=1cm]
  \item[Linter] strumento che analizza il codice per rilevare errori sintattici o stilistici.
  \item[Linguaggio di programmazione] mezzo formale per scrivere istruzioni comprensibili da un computer.
  \item[Login] procedura di autenticazione che consente l'accesso al sistema.
  \item[Log strutturato] registro di eventi con formato standard che facilita l'analisi automatizzata.
\end{description}

% M
\subsection*{M}
\begin{description}[leftmargin=2cm, labelindent=1cm]
  \item[Manutenibilità] facilità con cui il software può essere corretto o migliorato nel tempo.
  \item[MB] unità di misura della quantità di dati (megabyte).
  \item[Menù] elenco di opzioni tra cui l'utente può scegliere.
  \item[Messaggio di errore] testo mostrato quando si verifica un problema nell'esecuzione.
  \item[Metodologia] insieme di processi e strumenti adottati nello sviluppo software.
  \item[Metodologie di sviluppo] approcci sistematici come Agile, Waterfall.
  \item[Modale] finestra sovrapposta che richiede l'interazione dell'utente.
  \item[Modello di Cockburn] formalismo per la descrizione strutturata dei casi d'uso.
  \item[Modifiche strutturali] cambiamenti significativi nell'architettura o nella base di codice.
\end{description}

% N (nessuna voce)

% O
\subsection*{O}
\begin{description}[leftmargin=2cm, labelindent=1cm]
  \item[Requisiti operativi] requisiti che definiscono le condizioni di utilizzo e manutenzione del sistema.
\end{description}

% U
\subsection*{U}
\begin{description}[leftmargin=2cm, labelindent=1cm]
    \item[UI-(User Interface)] interfaccia grafica con cui l'utente interagisce.
    \item[URL] indirizzo che identifica una risorsa in rete.
    \item[URL param] parametro aggiunto a un URL per trasmettere informazioni.
    \item[Usabilità] facilità con cui l'utente può utilizzare un sistema.
    \item[Up time] tempo in cui il sistema risulta operativo e accessibile.
\end{description}

% V
\subsection*{V}
\begin{description}[leftmargin=2cm, labelindent=1cm]
    \item[Validazione] processo con cui si verifica che un sistema, prodotto o software soddisfi realmente i bisogni e le aspettative dell'utente finale.
    \item[Version Control] sistema che gestisce e traccia le modifiche apportate ai file di un progetto nel tempo.
\end{description}

% X
\subsection*{X}
\begin{description}[leftmargin=2cm, labelindent=1cm]
    \item[XSS] Vulnerabilità nelle applicazioni web che permette a un attaccante di inserire ed eseguire script non autorizzati nel browser di un altro utente
\end{description}